% GENERATED FILE. Edit canonical lessons, then run npm run generate:books.
% canonical-source-hash: fnv1a64:bd3d7fbbd0eccec8
% canonical-lessons: GU-C441-sarovar, GU-C441-dem, GU-C441-vav, GU-C441-katho, GU-C441-khupvu

\chapter{Lakes, Dams and Stepwells}
\label{ch:gu-lakes-dams-and-stepwells}
\hlchaptermodality{441}

\begin{chapteropening}
\textbf{By the end of this chapter:} \emph{I can talk about where a town gets its water.}

Complete the last lesson of chapter 441: I can talk about where a town gets its water.
\end{chapteropening}

\section[\texorpdfstring{\gu{સરોવર} (sarovar)}{સરોવર (sarovar)}]{\gu{સરોવર} (sarovar) — a large lake}

\label{lesson:GU-C441-sarovar}



\begin{quote}
Before the new one: say the Gujarati for a canal, then the Gujarati for to let flow.
\end{quote}

\subsection*{You'll want to know: \gu{સરોવર}}
\textbf{\gu{સરોવર}} — \emph{sarovar} — \textquotedblleft{}a large lake\textquotedblright{}.

\begin{quote}
\gu{આ} \gu{સરોવરમાં} \gu{ઘણી} \gu{માછલીઓ} \gu{છે}.
\end{quote}

\begin{quote}
\emph{ā sarovarmā̃ ghaṇī māchlīo che.}
\end{quote}

\begin{quote}
There are many fish in this lake.
\end{quote}

A neuter noun.

\begin{cousinweb}
\textbf{\gu{સરોવર}} (\emph{sarovar}) is a tatsama, Sanskrit \emph{sarovara}, \textquotedblleft{}a lake, a large pond\textquotedblright{} (Monier-Williams); a \gu{તળાવ} is usually smaller.
\end{cousinweb}

\begin{grammarlens}[title={What you've built}]
One more word for getting around.
\end{grammarlens}

\subsection*{Your turn}
\begin{itemize}
\raggedright
  \item \emph{Say it:} \emph{sarovar}
  \item \emph{Say it:} \emph{sarovar}, once more
  \item \emph{Say it:} \emph{ā sarovarmā̃ ghaṇī māchlīo che.} — There are many fish in this lake.
  \item \emph{Recall:} say the Gujarati for a canal, then the Gujarati for to let flow, then say \emph{sarovar} again
\end{itemize}

\begin{tcolorbox}[breakable,colback=teal!4,colframe=teal!35!black,title={Before you move on}]
What does \gu{સરોવર} mean? (\textquotedblleft{}A large lake\textquotedblright{}.) Say it once more.
\end{tcolorbox}
\section[\texorpdfstring{\gu{ડેમ} (ḍem)}{ડેમ (ḍem)}]{\gu{ડેમ} (ḍem) — a dam}

\label{lesson:GU-C441-dem}



\begin{quote}
Before the new one: say the Gujarati for to let flow, then the Gujarati for a large lake.
\end{quote}

\subsection*{You'll want to know: \gu{ડેમ}}
\textbf{\gu{ડેમ}} — \emph{ḍem} — \textquotedblleft{}a dam\textquotedblright{}.

\begin{quote}
\gu{નદી} \gu{પર} \gu{નવો} \gu{ડેમ} \gu{બન્યો} \gu{છે}.
\end{quote}

\begin{quote}
\emph{nadī par navo ḍem banyo che.}
\end{quote}

\begin{quote}
A new dam has been built on the river.
\end{quote}

A masculine noun.

\begin{cousinweb}
\textbf{\gu{ડેમ}} (\emph{ḍem}) is borrowed from English \emph{dam}.
\end{cousinweb}

\begin{grammarlens}[title={What you've built}]
One more word for getting around.
\end{grammarlens}

\subsection*{Your turn}
\begin{itemize}
\raggedright
  \item \emph{Say it:} \emph{ḍem}
  \item \emph{Say it:} \emph{ḍem}, once more
  \item \emph{Say it:} \emph{nadī par navo ḍem banyo che.} — A new dam has been built on the river.
  \item \emph{Recall:} say the Gujarati for to let flow, then the Gujarati for a large lake, then say \emph{ḍem} again
\end{itemize}

\begin{tcolorbox}[breakable,colback=teal!4,colframe=teal!35!black,title={Before you move on}]
What does \gu{ડેમ} mean? (\textquotedblleft{}A dam\textquotedblright{}.) Say it once more.
\end{tcolorbox}
\section[\texorpdfstring{\gu{વાવ} (vāv)}{વાવ (vāv)}]{\gu{વાવ} (vāv) — a stepwell}

\label{lesson:GU-C441-vav}



\begin{quote}
Before the new one: say the Gujarati for a large lake, then the Gujarati for a dam.
\end{quote}

\subsection*{You'll want to know: \gu{વાવ}}
\textbf{\gu{વાવ}} — \emph{vāv} — \textquotedblleft{}a stepwell\textquotedblright{}.

\begin{quote}
\gu{ગામની} \gu{જૂની} \gu{વાવમાં} \gu{હજી} \gu{પાણી} \gu{છે}.
\end{quote}

\begin{quote}
\emph{gāmnī jūnī vāvmā̃ hajī pāṇī che.}
\end{quote}

\begin{quote}
There is still water in the village's old stepwell.
\end{quote}

A feminine noun.

\begin{cousinweb}
\textbf{\gu{વાવ}} (\emph{vāv}) is a tadbhava of Sanskrit \emph{vāpī}, \textquotedblleft{}an oblong pond\textquotedblright{}; CDIAL lists the Gujarati word, \textquotedblleft{}a large well with steps\textquotedblright{} (CDIAL 11529): steps that go down to the water.
\end{cousinweb}

\begin{grammarlens}[title={What you've built}]
One more word for getting around.
\end{grammarlens}

\subsection*{Your turn}
\begin{itemize}
\raggedright
  \item \emph{Say it:} \emph{vāv}
  \item \emph{Say it:} \emph{vāv}, once more
  \item \emph{Say it:} \emph{gāmnī jūnī vāvmā̃ hajī pāṇī che.} — There is still water in the village's old stepwell.
  \item \emph{Recall:} say the Gujarati for a large lake, then the Gujarati for a dam, then say \emph{vāv} again
\end{itemize}

\begin{tcolorbox}[breakable,colback=teal!4,colframe=teal!35!black,title={Before you move on}]
What does \gu{વાવ} mean? (\textquotedblleft{}A stepwell\textquotedblright{}.) Say it once more.
\end{tcolorbox}
\section[\texorpdfstring{\gu{કાંઠો} (kā̃ṭho)}{કાંઠો (kā̃ṭho)}]{\gu{કાંઠો} (kā̃ṭho) — a riverbank; the bank of a lake}

\label{lesson:GU-C441-katho}



\begin{quote}
Before the new one: say the Gujarati for a dam, then the Gujarati for a stepwell.
\end{quote}

\subsection*{You'll want to know: \gu{કાંઠો}}
\textbf{\gu{કાંઠો}} — \emph{kā̃ṭho} — \textquotedblleft{}a riverbank; the bank of a lake\textquotedblright{}.

\begin{quote}
\gu{નદીના} \gu{કાંઠે} \gu{બાળકો} \gu{રમે} \gu{છે}.
\end{quote}

\begin{quote}
\emph{nadīnā kā̃ṭhe bāḷko rame che.}
\end{quote}

\begin{quote}
Children play on the river bank.
\end{quote}

A masculine noun.

\begin{cousinweb}
\textbf{\gu{કાંઠો}} (\emph{kā̃ṭho}) is a tadbhava of Sanskrit \emph{kaṇṭha}, \textquotedblleft{}the throat\textquotedblright{}, and so a rim or edge; CDIAL lists the Gujarati word, \textquotedblleft{}a bank, a coast\textquotedblright{} (CDIAL 2680).
\end{cousinweb}

\begin{grammarlens}[title={What you've built}]
One more word for getting around.
\end{grammarlens}

\subsection*{Your turn}
\begin{itemize}
\raggedright
  \item \emph{Say it:} \emph{kā̃ṭho}
  \item \emph{Say it:} \emph{kā̃ṭho}, once more
  \item \emph{Say it:} \emph{nadīnā kā̃ṭhe bāḷko rame che.} — Children play on the river bank.
  \item \emph{Recall:} say the Gujarati for a dam, then the Gujarati for a stepwell, then say \emph{kā̃ṭho} again
\end{itemize}

\begin{tcolorbox}[breakable,colback=teal!4,colframe=teal!35!black,title={Before you move on}]
What does \gu{કાંઠો} mean? (\textquotedblleft{}A riverbank; the bank of a lake\textquotedblright{}.) Say it once more.
\end{tcolorbox}
\section[\texorpdfstring{\gu{ખૂંપવું} (khū̃pvũ)}{ખૂંપવું (khū̃pvũ)}]{\gu{ખૂંપવું} (khū̃pvũ) — to sink in (mud); to get stuck deep}

\label{lesson:GU-C441-khupvu}



\begin{quote}
Before the new one: say the Gujarati for a stepwell, then the Gujarati for a riverbank.
\end{quote}

\subsection*{You'll want to know: \gu{ખૂંપવું}}
\textbf{\gu{ખૂંપવું}} — \emph{khū̃pvũ} — \textquotedblleft{}to sink in (mud); to get stuck deep\textquotedblright{}.

\begin{quote}
\gu{કાદવમાં} \gu{મારા} \gu{પગ} \gu{ખૂંપી} \gu{ગયા}.
\end{quote}

\begin{quote}
\emph{kādavmā̃ mārā pag khū̃pī gayā.}
\end{quote}

\begin{quote}
My feet sank into the mud.
\end{quote}

Past (it …): \gu{તે} \gu{ગાડું} \gu{કાદવમાં} \gu{ખૂંપી} \gu{ગયું}, \emph{te gāḍũ kādavmā̃ khū̃pī gayũ}.

\begin{cousinweb}
\textbf{\gu{ખૂંપવું}} (\emph{khū̃pvũ}) goes back to \emph{*skumpati}, \textquotedblleft{}presses in\textquotedblright{}, a form CDIAL reconstructs; CDIAL lists the Gujarati verb, \textquotedblleft{}to pierce\textquotedblright{} (CDIAL 13658).
\end{cousinweb}

\begin{grammarlens}[title={What you've built}]
One more word for getting around.
\end{grammarlens}

\subsection*{Your turn}
\begin{itemize}
\raggedright
  \item \emph{Say it:} \emph{khū̃pvũ}
  \item \emph{Say it:} \emph{khū̃pvũ}, once more
  \item \emph{Say it:} \emph{kādavmā̃ mārā pag khū̃pī gayā.} — My feet sank into the mud.
  \item \emph{Recall:} say the Gujarati for a stepwell, then the Gujarati for a riverbank, then say \emph{khū̃pvũ} again
\end{itemize}

\begin{tcolorbox}[breakable,colback=teal!4,colframe=teal!35!black,title={Before you move on}]
What does \gu{ખૂંપવું} mean? (\textquotedblleft{}To sink in (mud); to get stuck deep\textquotedblright{}.) Say it once more.
\end{tcolorbox}
